\section{Research Questions}

\begin{description}[style=nextline,leftmargin=0pt]
  \item[RQ1: Semantic boundary.]
  What canonical state projection is sufficient for observational equivalence
  with v1.02h gameplay across all reachable states in the evaluation corpus?

  \item[RQ2: Differential diagnosis.]
  How can heterogeneous executions---the original 32-bit Windows program and a
  portable Linux reconstruction---emit a stable, versioned representation that
  identifies the first divergent subsystem and field?

  \item[RQ3: Untrusted replay witnesses.]
  How can the compact input information in \code{.rpy} files be reused while
  preventing their per-stage checkpoints from becoming trusted state resets?

  \item[RQ4: Semantics-preserving slicing.]
  Which rendering, audio, animation, message, and effect operations can be
  removed, and which must remain as dependency-preserving shadow semantics due
  to timing, control flow, collision geometry, or shared RNG consumption?

  \item[RQ5: Arithmetic and slicing soundness.]
  Can the shipped program's x87, binary32 storage, and transcendental behavior
  be represented exactly at practical proof cost; and what machine-checked
  refinement justifies removing state or substituting integer arithmetic?

  \item[RQ6: Proof decomposition.]
  For each explicitly bounded slice, and after full semantic equivalence and
  refinement soundness for broader claims, which state layout, arithmetic
  model, frame chunking, and recursive aggregation strategy makes a spell,
  stage, and full six-stage run practical in a zkVM?
\end{description}

RQ1--RQ5 remain prerequisites for a whole-game answer to RQ6.  The OpenVM
experiments select a backend only for a retail-anchored Player slice.
The first proof hash-binds per-frame environment records; the second derives
those records inside a fail-closed enclosing player-state profile; and the
third derives focus, previous input, fire-timer cadence, and the
\code{SpawnBullets} request schedule.  All commit only fixed configuration,
replay masks, and their reached endpoints after the first experiment.  None
preselects the final whole-game architecture.  Corpus agreement may motivate a
projection or arithmetic design, but cannot by itself discharge RQ5's
universal claim.
